% ----------------------------------------------------------------------
%  Итоги главы 1: сводка, ответы, задачи, литература
% ----------------------------------------------------------------------
\clearpage
\section*{Итоги главы}
\addcontentsline{toc}{section}{Итоги главы}
\markright{Итоги главы 1}

Глава прошла путь от внешней формы кристалла к его внутренней геометрии
и от неё~— к свойствам и превращениям. Ключевая мысль: \vazhno{кристалл~—
это решётка плюс базис}, а почти всё остальное (индексы, симметрия,
число независимых констант, системы скольжения, ориентационные
соотношения) следует из геометрии решётки и её симметрии.
Таблица~\ref{tab:ch1-summary} сводит вместе инструменты главы.

\medskip\noindent\begin{minipage}{\textwidth}
\centering
\captionof{table}{Инструменты физической кристаллографии и их применение}
\label{tab:ch1-summary}
\footnotesize
\begin{tabularx}{\textwidth}{@{}>{\raggedright\arraybackslash\hspace{0pt}}p{30mm}>{\raggedright\arraybackslash\hspace{0pt}}X>{\raggedright\arraybackslash\hspace{0pt}}p{40mm}@{}}
\toprule
Инструмент & Что даёт & Где используется \\
\midrule
Стереографическая проекция, сетка Вульфа & углы между направлениями и плоскостями на плоском рисунке & ориентировка монокристаллов, полюсные фигуры, системы скольжения \\
Индексы Миллера (Миллера--Браве), обратная решётка & описание направлений и плоскостей; $d_{hkl}=1/|\vb H_{hkl}|$, закон зон & дифракция, зоны Бриллюэна, все расчёты углов \\
Матрицы поворота & смена систем координат, преобразование индексов & оси симметрии, разориентация зёрен, текстуры \\
Точечные и пространственные группы & 32 класса, 14 решёток Браве, 230 групп & классификация кристаллов, правила отбора свойств \\
Принципы Неймана и Кюри & симметрия свойств и кристалла во внешнем поле & пьезо- и пироэффект, анизотропия \\
Тензоры свойств & число независимых констант по сингониям & проводимость, упругость ($c_{ij}$), тепловое расширение \\
Фактор Шмида, тензор дисторсии & какая система скольжения работает и как поворачивается решётка & пластичность монокристаллов, текстуры \\
РСУ, $\Sigma$, ПРН & строение и энергия границ зёрен & зернограничная инженерия, диффузия по границам \\
Деформация Бейна, $P_1=RPB$ & решётка и форма мартенсита & закалка сталей, сплавы с памятью формы \\
\bottomrule
\end{tabularx}
\end{minipage}

\begin{itogi}[Главные формулы главы]
\begin{itemize}
\item Стереографическая проекция: $\rho=r\operatorname{tg}(\theta/2)$
\eqref{eq:stereo-rho}.
\item Закон зон: $hu+kv+lw=0$; межплоскостное расстояние $d_{hkl}=1/|\vb H_{hkl}|$,
для кубической решётки $d=a/\sqrt{h^2+k^2+l^2}$.
\item Угол между направлениями в кубическом кристалле
$\cos\varphi=\dfrac{u_1u_2+v_1v_2+w_1w_2}{\sqrt{u_1^2+v_1^2+w_1^2}\sqrt{u_2^2+v_2^2+w_2^2}}$
\eqref{eq:cos-cubic}.
\item Четырёхиндексные обозначения: $i=-(h+k)$; $[UVW]\to[uvtw]$ по
\eqref{eq:uvtw}.
\item Угол поворота по следу матрицы: $\operatorname{tr}R=\pm(1+2\cos\varphi)$.
\item Принцип Неймана $G_\text{свойства}\supseteq G_\text{кристалла}$;
принцип Кюри $G_\text{кр+возд}=G_\text{кр}\cap G_\text{возд}$.
\item Закон Гука $\sigma_{ij}=c_{ijkl}\varepsilon_{kl}$; фактор анизотропии
$A=2c_{44}/(c_{11}-c_{12})$ \eqref{eq:zener}; модуль Юнга кубического кристалла
\eqref{eq:E-cubic}.
\item Закон Шмида $\tau=\sigma\cos\chi\cos\lambda$ \eqref{eq:schmid};
тензор пластической дисторсии $u^0_{ij}=\alpha m_ih_j$ \eqref{eq:plastic-distortion}.
\item Двойникование: $s=2\operatorname{ctg}2\phi$ \eqref{eq:twin-shear}.
\item Малоугловая граница $D=b/\theta$; формула Франка \eqref{eq:frank};
энергия $E=E_0\theta(A-\ln\theta)$; критерий Брэндона
$\Delta\theta=15^\circ/\sqrt\Sigma$.
\item Деформация с инвариантной плоскостью $P=I+\vb d\tilde{\vb p}$;
деформация Бейна $\eta_1=\sqrt2a_\alpha/a_\gamma$, $\eta_3=c_\alpha/a_\gamma$.
\end{itemize}
\end{itogi}

%=======================================================================
\section*{Ответы на вопросы для самопроверки}
\addcontentsline{toc}{section}{Ответы на вопросы для самопроверки}
\markright{Ответы на вопросы главы 1}

{\small
\paragraph{1.1 (кристаллическое состояние)}
\begin{enumerate}[label=\arabic*.]
\item Форма (габитус) зависит от условий роста: скорости роста разных
граней различны, и одни грани развиваются сильнее других. Но все грани
параллельны одним и тем же плоскостям решётки, поэтому углы между
соответственными гранями одинаковы (закон постоянства углов).
\item Отношения параметров граней~— отношения небольших целых чисел;
значит, грани проходят через узлы трёхмерной периодической решётки:
кристалл имеет решётчатое строение.
\item Однородность~— одинаковость свойств в параллельно расположенных
точках (симметрия относительно трансляций); изотропия~— одинаковость
свойств по всем направлениям. Кристалл однороден, но анизотропен
(например, модуль Юнга меди вдоль $\nfam{111}$ почти втрое больше, чем
вдоль $\nfam{100}$).
\item Решётка~— геометрическое множество эквивалентных точек; её начало
можно поместить в любую точку структуры. Атомы базиса располагаются
относительно узлов, и в общем случае ни один из них не обязан попадать
в узел.
\item Три атома (A, B, C).
\end{enumerate}

\paragraph{1.2 (основы кристаллографии)}
\begin{enumerate}[label=\arabic*.]
\item Стереографическая проекция плоскости~— её след (дуга большого
круга), гномостереографическая~— полюс (проекция нормали). В
кубическом кристалле направление $[hkl]$ перпендикулярно плоскости
$(hkl)$, поэтому полюс $(hkl)$ совпадает с выходом $[hkl]$.
\item Угол между $[001]$ и нормалью к $(011)$ равен $45^\circ$:
$\rho=r\operatorname{tg}22{,}5^\circ=4{,}14$~см.
\item Обратные отрезки $\frac12:\frac13:(-1)$, умноженные на 6: $(32\br6)$.
\item $[210]\to[10\br10]$, $[1\br10]\to[1\br100]$.
\item Векторное произведение нормалей: $[11\br2]$.
\item $1+2\cos\varphi=-1$, $\varphi=180^\circ$ (ось 2).
\item $d_{110}=a/\sqrt2=0{,}2027$~нм, $d_{111}=a/\sqrt3=0{,}1655$~нм. (В
ОЦК-решётке посередине между плоскостями $(111)$ лежат такие же атомные
плоскости, поэтому реальное расстояние между атомными слоями $\parallel(111)$~—
$d_{222}=0{,}083$~нм; отражение $111$ погасает.)
\item Вторые соседи в ОЦК-решётке (на расстоянии $a$) лишь на 15\,\%
дальше первых ($a\sqrt3/2$), и плоскости $\plfam{100}$, делящие пополам
отрезки до них, срезают вершины октаэдра из плоскостей $\plfam{111}$:
получается усечённый октаэдр (8 шестиугольников и 6 квадратов).
\end{enumerate}

\paragraph{1.3 (симметрия кристаллов)}
\begin{enumerate}[label=\arabic*.]
\item Поворот на $72^\circ$ несовместим с трансляционной периодичностью:
пятиугольниками нельзя без зазоров замостить плоскость. У квазикристаллов
трансляционной периодичности нет (порядок квазипериодический), поэтому
запрет на них не распространяется.
\item Ось 6: $1+2\cos60^\circ=2$; ось $\br4$: $-(1+2\cos90^\circ)=-1$.
\item $\br2=iC_2$: при оси 2 вдоль $z$ матрица $\operatorname{diag}(-1,-1,1)$,
умноженная на $-I$, даёт $\operatorname{diag}(1,1,-1)$~— отражение в
плоскости $\perp z$. $\br6=iC_6$: $(iC_6)^2=C_3$, $(iC_6)^3=iC_2=m_\perp$,
т.\,е. $\br6=3/m$.
\item 48. Подгрупп порядка 5 нет: по теореме Лагранжа порядок подгруппы
делит порядок группы, а 5 не делит 48.
\item Класс $m\br3m$ центросимметричен~— пьезоэффекта нет. Класс $\br43m$
нецентросимметричен и пьезоэлектричен (сфалерит, GaAs).
\item Тетрагональная $C$-решётка~— та же примитивная тетрагональная
решётка $P$ с ячейкой в два раза меньшего объёма, повёрнутой на $45^\circ$
вокруг оси $c$ (рёбра основания $a/\sqrt2$).
\item За поворот на $90^\circ$ ось $4_1$ сдвигает на $c/4$, ось $4_3$~— на
$3c/4$ (эквивалентно $-c/4$): это правый и левый винты, зеркальные
двойники друг друга.
\item Магнитное поле~— аксиальный вектор: $\infty/m$ (вращающийся цилиндр:
центр симметрии и плоскость $\perp\vb H$, но нет продольных плоскостей).
Электрическое поле~— полярный вектор: $\infty mm$ (конус).
\end{enumerate}

\paragraph{1.4 (структура кристаллов)}
\begin{enumerate}[label=\arabic*.]
\item Обе~— плотнейшие упаковки одинаковых плотных слоёв, отличающиеся
лишь порядком чередования слоёв (ABC и AB): плотность заполнения
одинакова, 0,74.
\item ГЦК: одна октаэдрическая и две тетраэдрические пустоты на атом.
ОЦК: три октаэдрические (сплюснутые) и шесть тетраэдрических.
\item В аустените октапустоты крупные ($r=0{,}414R$). В феррите атом
углерода в сплюснутом октаэдре раздвигает лишь два ближайших атома вдоль
оси сплющивания; в тетрапустоте пришлось бы раздвигать четыре атома, что
энергетически невыгодно.
\item Половину (через одну); у каждого атома четыре ближайших соседа в
вершинах тетраэдра~— направленные $sp^3$-связи.
\item Обе построены из тетраэдров; в сфалерите слои упакованы по закону
ABC (кубическая решётка, $F\br43m$), в вюрците~— AB (гексагональная,
$P6_3mc$).
\item Атомы в вершинах (Cs) и в центре (Cl) разные, поэтому трансляции
$\frac12\frac12\frac12$ нет: решётка простая кубическая с базисом из двух
атомов.
\item 12 атомов Fe и 4 атома Al (16 атомов, Fe$_3$Al).
\item Поверхностная реконструкция грани $(111)$ кремния: двумерная ячейка
поверхности в $7\times7$ раз больше ячейки идеальной (нереконструированной)
грани.
\end{enumerate}

\paragraph{1.5 (физические свойства)}
\begin{enumerate}[label=\arabic*.]
\item По принципу Кюри $m\br3m\cap\infty/mmm$ (ось цилиндра $\parallel[110]$)
$=mmm$: остаются оси 2 вдоль $[110]$, $[\br110]$, $[001]$ и
перпендикулярные им плоскости (ромбическая симметрия).
\item Пироэлектрик должен иметь единственное полярное направление, а в
кубических группах четыре оси 3 вдоль $\nfam{111}$ не допускают его.
\item Проводимость~— тензор 2-го ранга, а у кубического кристалла все
такие свойства изотропны; модуль Юнга определяется тензором 4-го ранга,
у которого три независимые компоненты, и анизотропия сохраняется.
\item $E_{\nfam{110}}\approx130$~ГПа (для сравнения $E_{\nfam{100}}=67$,
$E_{\nfam{111}}=190$~ГПа).
\item Два ($\sigma_{11}=\sigma_{22}$ и $\sigma_{33}$): в~\eqref{eq:n-invariants}
вклады $E$, $C_2$, $i$, $\sigma_h$ равны 6, 2, 6, 2, а у $C_4$, $C_4^3$, $S_4$,
$S_4^3$~— нули; $n=16/8=2$, как у всех тетрагональных классов.
\item При $c/a=1{,}86$ базисные слои цинка далеко отстоят друг от друга и
связаны слабо; поэтому кристалл легко сжимается вдоль оси $c$.
\end{enumerate}

\paragraph{1.6 (пластическая деформация)}
\begin{enumerate}[label=\arabic*.]
\item Скольжение сдвигает на вектор трансляции, и узлы попадают в узлы той
же решётки. При двойниковании атомы смещаются на доли трансляции, и
решётка переходит в зеркально-симметричную ориентацию.
\item Когда $\vb n$, $\vb m$ и $F$ лежат в одной плоскости и
$\chi=\lambda=45^\circ$. Больше $\frac12$ быть не может: $\lambda\ge90^\circ-\chi$,
поэтому $\cos\chi\cos\lambda\le\cos\chi\sin\chi\le\frac12$.
\item $\cos\chi=1/\sqrt3$, $\cos\lambda=1/\sqrt2$, фактор $0{,}408$. Такой
же фактор у 8 систем (в каждой из четырёх плоскостей $\plfam{111}$~— два
направления, не перпендикулярные $[001]$).
\item Базисное скольжение даёт лишь две независимые системы (меньше
пяти) и не может изменить размер кристалла вдоль оси $c$: зёрна не могут
согласованно деформироваться, на границах возникают трещины.
\item $\cos2\phi=\frac{1\cdot1+1\cdot1+2\cdot(-2)}{6}=-\frac13$, т.\,е. угол
между плоскостями $70{,}53^\circ$; $s=2\operatorname{ctg}70{,}53^\circ=0{,}707$.
\item Точка в центре проекции (полюс $(001)$) и основной круг (полюса
$(100)$ и $(010)$ при вращении вокруг $F$).
\end{enumerate}

\paragraph{1.7 (границы зёрен)}
\begin{enumerate}[label=\arabic*.]
\item Пять (три~— разворот, два~— нормаль к границе); на атомном уровне
ещё три (вектор $\vb t$), всего восемь.
\item У границы наклона ось разворота лежит в плоскости границы, у границы
кручения перпендикулярна ей. Малоугловая граница наклона~— стенка
краевых дислокаций, кручения~— сетка винтовых.
\item $\theta=b/D=0{,}256/50=5{,}1\cdot10^{-3}$~рад $\approx0{,}29^\circ$.
\item Поворот на $90^\circ$ вокруг $[001]$~— операция симметрии куба, а
$53{,}13^\circ=90^\circ-36{,}87^\circ$: это одна и та же разориентация,
описанная по-разному (и с переставленными зёрнами); совпадающие узлы те же.
\item Граница $\Sigma3$ по $\plfam{111}$~— зеркальная плоскость: у атомов
сохраняются все ближайшие соседи, меняется лишь порядок чередования
плотных слоёв (\ldots ABCACB\ldots); энергия порядка половины энергии
дефекта упаковки.
\item Нет: наибольший угол разориентации кубических кристаллов
$62{,}8^\circ$, и у любого поворота на $70^\circ$ есть эквивалентное описание
с меньшим углом (например, $70{,}5^\circ$ вокруг $\nfam{110}$~— это
$60^\circ$ вокруг $\nfam{111}$, двойник $\Sigma3$).
\end{enumerate}

\paragraph{1.8 (мартенситные превращения)}
\begin{enumerate}[label=\arabic*.]
\item Поверхность превращённой области остаётся плоской (деформация
однородна), а риски~— прямыми и непрерывными, изламываясь только на
следе габитуса (плоскость габитуса не искажается и не поворачивается).
\item Деформация решётки меняет кристаллическую структуру (ГЦК$\to$ОЦТ);
деформация с инвариантной решёткой меняет только форму, не меняя
решётки: скольжение, двойникование.
\item $\eta_1=\sqrt2\cdot0{,}2850/0{,}3584=1{,}125$, $\eta_3=0{,}2958/0{,}3584=0{,}825$;
$\Delta V/V=\eta_1^2\eta_3-1\approx+4{,}4\,\%$.
\item У неё два главных удлинения больше 1, одно меньше и ни одно не
равно 1; неискажённые направления образуют конус, но не плоскость.
\item $1\cdot1+0\cdot1+(-1)\cdot1=0$ и $0\cdot1+1\cdot1+1\cdot(-1)=0$: оба
условия зональности выполнены.
\end{enumerate}
}

%=======================================================================
\section*{Задачи}
\addcontentsline{toc}{section}{Задачи}
\markright{Задачи к главе 1}

Задачи взяты из раздела «Контрольные вопросы, задачи и упражнения»
исходного учебника (номера в скобках) и снабжены ответами.

\begin{enumerate}[label={\sffamily\bfseries\color{fmHead}\arabic*.}]
\item (13) Найдите индексы направлений, проходящих через узлы
$\uzel{000}$ и $\uzel{1\,1\,\br{\tfrac12}}$; $\uzel{\tfrac13\,1\,1}$ и
$\uzel{0\,\tfrac34\,1}$; $\uzel{\br{\tfrac23}\,0\,1}$ и
$\uzel{1\,0\,\tfrac12}$.\\
\emph{Ответ:} $[22\br1]$; $[430]$; $[10\,0\,\br3]$.
\item (15) Найдите индексы плоскостей, отсекающих на осях отрезки
$1,\,-\frac12,\,\frac13$; $-\frac13,\,1,\,\frac12$; $\frac12,\,\frac14,\,\frac12$.\\
\emph{Ответ:} $(1\br23)$; $(\br312)$; $(121)$.
\item (17) Входят ли плоскости $(12\br3)$, $(111)$, $(01\br1)$ в зону с осью
$[111]$?\\
\emph{Ответ:} $(12\br3)$ и $(01\br1)$~— да ($hu+kv+lw=0$), $(111)$~— нет.
\item (19) Найдите ось зоны, содержащей плоскости $(01\br11)$ и $(000\br1)$,
в трёх- и четырёхиндексной записи.\\
\emph{Ответ:} $[100]=[2\br1\br10]$.
\item (20) Найдите в кубическом кристалле углы между $[100]$ и $[110]$;
$[110]$ и $[11\br1]$; $[123]$ и $(123)$; $(100)$ и $(1\br10)$; $(100)$ и
$(1\br1\br1)$.\\
\emph{Ответ:} $45^\circ$; $35{,}26^\circ$; $0^\circ$; $45^\circ$; $54{,}74^\circ$.
\item (26) Найдите угол между базисными векторами примитивной ячейки
ГЦК-решётки.\\
\emph{Ответ:} $60^\circ$ (векторы $\frac a2[110]$, $\frac a2[011]$,
$\frac a2[101]$).
\item (28) Какие индексы будут у $(1\br10)$ и $[1\br10]$ кубического кристалла в
базисе примитивной ячейки ГЦК-решётки~\eqref{eq:prim}?\\
\emph{Ответ:} $h'=Ah$ даёт $(0\,\br{\tfrac12}\,\tfrac12)\to(0\br11)$; вектор
$a[1\br10]=2(\vb p_3-\vb p_2)$, т.\,е. $[0\br22]$, или направление $[0\br11]$.
\item (29--32) Укажите число и индексы граней ячейки Вигнера--Зейтца и
первой зоны Бриллюэна ГЦК- и ОЦК-решёток.\\
\emph{Ответ:} ячейка ВЗ ГЦК~— ромбододекаэдр, 12 граней $\plfam{110}$;
зона Бриллюэна ГЦК~— усечённый октаэдр, 8 граней $\plfam{111}$ и 6 граней
$\plfam{200}$; ячейка ВЗ ОЦК~— усечённый октаэдр, 8 граней $\plfam{111}$ и
6 граней $\plfam{100}$; зона Бриллюэна ОЦК~— ромбододекаэдр, 12 граней
$\plfam{110}$.
\item (34--36) Запишите матрицы поворотов: $C_4$ вокруг $[010]$ на
$90^\circ$ по часовой стрелке; $C_3$ вокруг $[\br111]$ на $120^\circ$
против часовой стрелки; $C_2$ вокруг $[110]$.\\
\emph{Ответ:}
$\begin{pmatrix}0&0&\br1\\0&1&0\\1&0&0\end{pmatrix}$,
$\begin{pmatrix}0&\br1&0\\0&0&1\\\br1&0&0\end{pmatrix}$,
$\begin{pmatrix}0&1&0\\1&0&0\\0&0&\br1\end{pmatrix}$.
\item (37) Найдите $C=C_2[01\br1]\,C_4[001]$ ($C_4$~— на $90^\circ$ против
часовой стрелки) и определите тип и индексы оси.\\
\emph{Ответ:} $C=\begin{pmatrix}0&1&0\\0&0&\br1\\\br1&0&0\end{pmatrix}$,
$\det C=1$, $\operatorname{tr}C=0$: поворотная ось 3 вдоль $[11\br1]$.
\item (38--41) Найдите результат последовательного выполнения
операций: $\sigma_{(1\br10)}\sigma_{(101)}$; $\sigma_{(011)}\,i$;
$C_3[\br111]\,i$; $C_2[101]\,\sigma_{(01\br1)}$.\\
\emph{Ответ:} поворотная ось 3 вдоль $[11\br1]$ (две плоскости под
углом $60^\circ$ дают поворот на $120^\circ$ вокруг линии их пересечения);
поворотная ось 2 вдоль $[011]$; инверсионная ось $\br3$ вдоль $[\br111]$;
инверсионная ось $\br3$ вдоль $[\br111]$ (определитель $-1$, угол
$120^\circ$).
\item (46) Почему в кубической сингонии нет базоцентрированной решётки
Браве?\\
\emph{Ответ:} центрирование одной пары граней несовместимо с осями 3
вдоль $\nfam{111}$, которые переводят грани $(001)$ в $(100)$ и $(010)$;
такая решётка имела бы лишь тетрагональную симметрию.
\item (49, 50, 55) Найдите координационные числа и радиусы первых пяти
координационных сфер в ГЦК, ОЦК-структурах и в алмазе.\\
\emph{Ответ:} ГЦК: $a\sqrt{n/2}$ ($n=1\ldots5$), числа 12, 6, 24, 12, 24.
ОЦК: $a\sqrt3/2$, $a$, $a\sqrt2$, $a\sqrt{11}/2$, $a\sqrt3$; числа 8, 6, 12,
24, 8. Алмаз: $\frac a4\sqrt3$, $\frac a4\sqrt8$, $\frac a4\sqrt{11}$, $a$,
$\frac a4\sqrt{19}$; числа 4, 12, 12, 6, 12.
\item (53) Докажите, что в идеальной ГПУ-структуре $c/a=\sqrt{8/3}=1{,}633$.\\
\emph{Ответ:} атом слоя B лежит над центром правильного треугольника из
атомов слоя A со стороной $a$ на расстоянии $a$ от каждого: высота
правильного тетраэдра $h=a\sqrt{2/3}$, а $c=2h$.
\item (57, 58) Найдите координационные числа и расстояния в структурах
UO$_2$ (флюорит) и CaTiO$_3$ (перовскит).\\
\emph{Ответ:} UO$_2$: U по O~— 8 на $a\sqrt3/4$, O по U~— 4 на
$a\sqrt3/4$. Перовскит (Ti в центре куба): Ca--Ti~— 8 на $a\sqrt3/2$,
Ti--O~— 6 на $a/2$, O--Ca~— 4 на $a/\sqrt2$, Ti--Ca~— 8 на $a\sqrt3/2$,
O--Ti~— 2 на $a/2$, Ca--O~— 12 на $a/\sqrt2$.
\item (59--61) Какие элементы симметрии сохранятся: у класса $\br43m$ при
растяжении вдоль $[010]$; у $m\br3m$ в электрическом поле вдоль $[1\br11]$; у
$6/mmm$ в магнитном поле вдоль $[01\br10]$?\\
\emph{Ответ:} $\br42m$ (ось $\br4$ вдоль $[010]$, оси 2 вдоль $[100]$ и
$[001]$, плоскости $(101)$ и $(10\br1)$); $3m$ (ось 3 вдоль $[1\br11]$ и
три плоскости $\plfam{110}$, проходящие через неё); $2/m$ (ось 2 вдоль
поля, плоскость $\perp$ полю, центр инверсии).
\item (62, 63) Найдите модули Юнга вдоль $\nfam{100}$, $\nfam{110}$,
$\nfam{111}$ монокристаллов Cu ($c_{11}=168$, $c_{12}=121$, $c_{44}=76$~ГПа)
и Nb ($c_{11}=246$, $c_{12}=134$, $c_{44}=29$~ГПа).\\
\emph{Ответ:} Cu: 67, 131, 192~ГПа ($A=3{,}2$, максимум вдоль $\nfam{111}$);
Nb: 152, 93, 82~ГПа ($A=0{,}52$, максимум вдоль $\nfam{100}$).
\item (64) То же для Be вдоль $[2\br1\br10]$, $[1\br100]$, $[\br1011]$ ($c_{11}=292$,
$c_{12}=27$, $c_{13}=14$, $c_{33}=336$, $c_{44}=162$~ГПа, $c/a=1{,}57$).\\
\emph{Ответ:} в базисной плоскости гексагональный кристалл упруго
изотропен: $E=289$~ГПа вдоль обоих направлений; вдоль $[\br1011]$
(угол $47{,}8^\circ$ с осью $c$) $E\approx321$~ГПа; вдоль $[0001]$~— 335~ГПа.
\item (65) ГЦК-кристалл растянут напряжением $\sigma$ вдоль $[\br132]$. Найдите
касательные напряжения в первичной, сопряжённой, критической и
поперечной системах.\\
\emph{Ответ:} ось лежит в треугольнике $[010]$--$[011]$--$[\br111]$;
первичная $(111)[\br110]$: $0{,}467\sigma$; сопряжённая $(\br11\br1)[011]$:
$0{,}292\sigma$; критическая $(\br111)[110]$: $0{,}350\sigma$; поперечная
$(11\br1)[\br110]$: 0 (ось лежит в этой плоскости).
\item (67) Первичная система ГЦК-кристалла $(111)[10\br1]$. Какова
поперечная система?\\
\emph{Ответ:} $(1\br11)[10\br1]$~— вторая плоскость $\plfam{111}$, содержащая
то же направление скольжения.
\item (70--73) Сколько независимых компонент у тензора пластической
дисторсии? Сколько систем скольжения $\plfam{112}\nfam{111}$ в ОЦК-металлах,
базисных и призматических систем в ГПУ-металлах?\\
\emph{Ответ:} 8 (след равен нулю: 5~— деформация, 3~— поворот);
12; 3; 3.
\item (78) Текстура прокатки $\plfam{101}\nfam{112}$. Каково поперечное
направление?\\
\emph{Ответ:} $\text{ПН}=\text{НН}\times\text{НП}$, например
$(101)\times[\br121]\parallel[\br1\br11]$, т.\,е. $\nfam{111}$.
\item (80) Найдите углы Эйлера $\theta$, $\psi$ стержня ФРО для аксиальной
текстуры $\nfam{mnp}$.\\
\emph{Ответ:} $\cos\theta=p/\sqrt{m^2+n^2+p^2}$, $\operatorname{tg}\psi=m/n$
\eqref{eq:fiber-euler}.
\item (81, 82) Найдите направление линии краевой дислокации: в
ГЦК-кристалле с $\vb b\parallel[110]$ в плоскости $\plfam{111}$; в ОЦК с
$\vb b\parallel[111]$ в плоскости $\plfam{112}$.\\
\emph{Ответ:} $\vb l\parallel\vb b\times\vb n$: в плоскости $(1\br11)$~—
$[1\br1\br2]$, в $(\br111)$~— $[1\br12]$ (тип $\nfam{112}$); в ОЦК в
плоскости $(11\br2)$~— $[\br110]$ (тип $\nfam{110}$).
\item (83) На какие частичные дислокации Шокли расщепляется полная
дислокация $\frac a2[1\br10]$ в плоскости $(111)$?\\
\emph{Ответ:} $\frac a2[1\br10]\to\frac a6[2\br1\br1]+\frac a6[1\br21]$.
\item (87, 88) Найдите расстояние между дислокациями в симметричной
границе наклона в Cu при $\theta=1{,}5^\circ$ ($a=0{,}362$~нм) и угол
разворота в W при расстоянии 2,5~мкм ($a=0{,}316$~нм).\\
\emph{Ответ:} Cu: $b=a/\sqrt2=0{,}256$~нм, $D=b/\theta=9{,}8$~нм;
W: $b=a\sqrt3/2=0{,}274$~нм, $\theta=1{,}1\cdot10^{-4}$~рад $\approx0{,}006^\circ$.
\item (89) Найдите плоскость малоугловой границы кручения из винтовых
дислокаций с $\vb b_1=\frac a2[101]$ и $\vb b_2=\frac a2[110]$.\\
\emph{Ответ:} обе дислокации лежат в границе: $\vb n\parallel\vb b_1\times\vb b_2$,
плоскость $(\br111)$.
\item (92) Найдите матрицу деформации Бейна при превращении аустенита
($a=0{,}3584$~нм) в мартенсит ($a=0{,}2850$, $c=0{,}2958$~нм).\\
\emph{Ответ:} $B=\operatorname{diag}(1{,}125;\ 1{,}125;\ 0{,}825)$.
\item (93) Найдите углы между направлениями $[101]$ и $[11\br1]$, $[011]$ и
$[\br101]$ аустенита ($a=0{,}3610$~нм) после превращения в мартенсит
($a=0{,}2845$, $c=0{,}3020$~нм; ось $c$ вдоль $[001]_\gamma$).\\
\emph{Ответ:} $\eta_1=1{,}115$, $\eta_3=0{,}837$; углы меняются с
$90^\circ$ на $77{,}4^\circ$ и с $60^\circ$ на $68{,}9^\circ$.
\end{enumerate}

\paragraph{Вопросы для повторения (по исходному учебнику)}
{\small
\begin{enumerate}[label=\arabic*.]
\item Что такое кристалл в современном понимании?
\item Что такое габитус кристалла, соответственные плоскости кристалла?
\item Сформулируйте закон постоянства углов и закон рациональных
параметров.
\item Какие свойства кристалла являются макроскопическими, какие~—
микроскопическими?
\item Что такое гомологические точки кристалла? Сформулируйте
математические постулаты кристаллического пространства.
\item Что такое шар однородности и шар дискретности?
\item В чём отличие решётки и структуры кристалла?
\item На стереографической проекции показана зона плоскостей. Как будет
выглядеть та же зона на гномостереографической проекции?
\item Что такое кристаллографические индексы направления и плоскости?
\item Покажите в элементарной ячейке плоскости $(321)$, $(\br120)$; в
гексагональной ячейке~— плоскости $(1\br102)$, $(11\br20)$ и направления
$[1\br102]$, $[11\br20]$.
\item Совпадает ли точка первого рода $(011)$ с точкой второго рода $[011]$
в тетрагональном кристалле? Докажите, что в гексагональном кристалле
$\plfam{hki0}$ и $\nfam{hki0}$ совпадают.
\item Почему в кристаллах нет поворотных осей пятого порядка?
\item Перечислите групповые постулаты. Что такое генератор группы? Найдите
обратный элемент для $a^t$ циклической группы порядка $n$.
\item Какой физический смысл имеет вхождение элементов симметрии в один
класс сопряжённых элементов?
\item Покажите плоскость скольжения $n$-типа в ОЦК-решётке, $c$-типа в
ГЦК-решётке, $d$-типа и винтовую ось $4_1$ в алмазе.
\item Какая решётка Браве у ГПУ-структуры $\alpha$-Zr? Обладает ли
сфалерит центром симметрии?
\item Изобразите ППФ $(211)$ для текстуры $F\parallel\nfam{110}$; ОПФ направления
$F$ и направления $G\perp F$ для текстуры $\nfam{112}$; ФРО для текстуры
$\nfam{112}$.
\item Найдите выражения для касательного напряжения в системе
$(1\br1\br1)[10\br1]$ кубического кристалла при двухосном нагружении ($\sigma_1$
вдоль $[100]$, $\sigma_2$ вдоль $[001]$).
\end{enumerate}
}

%=======================================================================
\section*{Литература}
\addcontentsline{toc}{section}{Литература к главе 1}
\markright{Литература к главе 1}

{\small
\paragraph{Использованная в исходном учебнике}
\begin{enumerate}[label=\arabic*.]
\item Вайнштейн Б.\,К. Современная кристаллография. Т.~1. Симметрия
кристаллов. Методы структурной кристаллографии.~— М.: Наука, 1979.~— 384~с.
\item Вайнштейн Б.\,К., Фридкин В.\,М., Инденбом В.\,Л. Современная
кристаллография. Т.~2. Структура кристаллов.~— М.: Наука, 1979.~— 312~с.
\item Шувалов Л.\,А., Урусовская А.\,А., Желудев И.\,С. и др. Современная
кристаллография. Т.~4. Физические свойства кристаллов.~— М.: Наука,
1981.~— 496~с.
\item Кристиан Дж. Теория превращений в металлах и сплавах. Ч.~1.~— М.:
Мир, 1978.~— 806~с.
\item Васильев Д.\,М. Физическая кристаллография.~— М.: Металлургия,
1972.~— 280~с.; 2-е изд.~— 1981.~— 256~с.
\item Русаков А.\,А. Рентгенография металлов.~— М.: Атомиздат, 1977.~— 480~с.
\item Русаков А.\,А., Яльцев В.\,Н. Основы рентгенографии металлов. Ч.~I.~—
М.: МИФИ, 1994.~— 172~с.
\item Чупрунов Е.\,В., Хохлов А.\,Ф., Фаддеев М.\,А. Основы
кристаллографии.~— М.: Физматлит, 2004.~— 500~с.
\item Уманский Я.\,С., Скаков Ю.\,А., Иванов А.\,Н., Расторгуев Л.\,Н.
Кристаллография, рентгенография и электронная микроскопия.~— М.:
Металлургия, 1982.~— 632~с.
\item Шаскольская М.\,П. Кристаллография.~— М.: Высшая школа, 1976.~— 391~с.
\item Штремель М.\,А. Прочность сплавов. Ч.~I. Дефекты решётки.~— М.:
Металлургия, 1982.~— 280~с.; Ч.~II. Деформация.~— М.: МИСиС, 1997.~— 527~с.
\item Новиков И.\,И., Розин К.\,М. Кристаллография и дефекты
кристаллической решётки.~— М.: Металлургия, 1990.~— 336~с.
\item Орлов А.\,Н., Перевезенцев В.\,Н., Рыбин В.\,В. Границы зёрен в
металлах.~— М.: Металлургия, 1980.~— 156~с.
\item Келли А., Гровс Г. Кристаллография и дефекты в кристаллах.~— М.:
Мир, 1974.~— 496~с.
\item Физическое металловедение / под ред. Р.\,У.~Кана. Вып.~1--3.~— М.:
Мир, 1967--1968.
\item Физическое металловедение / под ред. Р.\,У.~Кана, П.~Хаазена. 3-е
изд. В 3~т.~— М.: Металлургия, 1987.
\item Пенкаля Т. Очерки кристаллохимии.~— Л.: Химия, 1974.~— 496~с.
\item Баррет Ч.\,С., Масальский Т.\,Б. Структура металлов. В 2~ч.~— М.:
Металлургия, 1984.~— 352~с.
\item Вустер У. Применение тензоров и теории групп для описания
физических свойств кристаллов.~— М.: Мир, 1977.~— 384~с.
\end{enumerate}

\paragraph{Дополнительно}
\begin{enumerate}[label=\arabic*., resume]
\item Най Дж. Физические свойства кристаллов и их описание при помощи
тензоров и матриц.~— М.: Мир, 1967.
\item Сиротин Ю.\,И., Шаскольская М.\,П. Основы кристаллофизики.~— М.:
Наука, 1979.
\item Вассерман Г., Гревен И. Текстуры металлических материалов.~— М.:
Металлургия, 1969.
\item Кайбышев О.\,А., Валиев Р.\,З. Границы зёрен и свойства металлов.~—
М.: Металлургия, 1987.
\item Курдюмов Г.\,В., Утевский Л.\,М., Энтин Р.\,И. Превращения в железе и
стали.~— М.: Наука, 1977.
\end{enumerate}
}
